\documentclass[12pt]{article}

\usepackage[utf8]{inputenc}
\usepackage[T1]{fontenc}
\usepackage[brazil]{babel}

\usepackage{geometry}
\geometry{margin=2.5cm}
\usepackage{setspace}
\onehalfspacing

\usepackage{amsmath, amssymb, amsthm}
\usepackage{mathtools}

\numberwithin{equation}{section}

\theoremstyle{plain}
\newtheorem{theorem}{Teorema}[section]
\newtheorem{proposition}[theorem]{Proposição}
\newtheorem{lemma}[theorem]{Lema}
\newtheorem{corollary}[theorem]{Corolário}

\theoremstyle{definition}
\newtheorem{definition}[theorem]{Definição}
\newtheorem{example}[theorem]{Exemplo}
\newtheorem{exercise}{Exercício}[section]

\theoremstyle{remark}
\newtheorem*{remark}{Observação}

\usepackage{hyperref}
\usepackage{csquotes}

\usepackage[
    backend=biber,
    style=authoryear,
    maxcitenames=2
]{biblatex}

\addbibresource{/mnt/c/Users/nicho/Documents/nicholasvoltani.github.io/bibliography.bib}

\newcommand{\R}{\mathbb{R}}
\newcommand{\pref}{\succeq}
\newcommand{\strictpref}{\succ}
\newcommand{\indiff}{\sim}

% ------------------------------------------------
% Title (fill manually)
% ------------------------------------------------
\title{}
\author{}
\date{}

\begin{document}

\maketitle

\begin{longtable}[]{@{}l@{}}
\toprule\noalign{}
\endhead
\bottomrule\noalign{}
\endlastfoot
date: ``2026-03-30'' \\
tags: \\
- economics \\
aliases: \\
\end{longtable}

\hypertarget{conceitos-introdutuxf3rios}{%
\section{Conceitos introdutórios}\label{conceitos-introdutuxf3rios}}

\hypertarget{loterias}{%
\subsection{Loterias}\label{loterias}}

\begin{definition}

\end{definition}

\hypertarget{teoria-da-firma}{%
\subsection{Teoria da firma}\label{teoria-da-firma}}

\begin{definition}
O equivalente certo de um tomador de risco é a quantidade monetária
não-aleatória com a qual ele obtém utilidade igual à sua
\textit{utilidade média}.

Dada uma loteria \(L\) e uma utilidade de Bernoulli \(u(\cdot)\), o
equivalente certo (ou equivalente de certeza) \(c(L, u)\) é o valor
(monetário) \textit{com probabilidade \(100\%\) de ser recebido} que faz
com que a utilidade de Bernoulli se iguale à Utilidade de von
Neumann-Morgenstern \(U(\cdot)\) da loteria \(L\),
i.e.~\(U(L)\).\footnote{\textcite{Cowell2004}, p.~188 comenta sobre a
  confusão que ambas estas utilidades --- utilidade ``de Bernoulli''
  (cf. \textcite{Mas-Colell1995}, p.~184) e utilidade de von
  Neumann-Morgenstern --- podem suscitar: ``Here we encounter a
  terminologically awkward corner. We should not really call \(u\)
  {[}''utilidade de Bernoulli''{]} `the utility function' because the
  whole expression {[}utilidade de von Neumann-Morgenstern,
  \(U(L) = \sum \limits_{k} p_{k} u_{k}\){]} is the person's utility; so
  \(u\) is sometimes known as the individual's \textit{cardinal utility
  function} or \textit{felicity function}; arguably neither term is a
  particularly happy choice of words.''.}

Formalmente, temos que, para uma loteria \((p_{1}, \dots, p_{n})\)
relacionada às utilidades (de Bernoulli)
\((u_{1} \coloneqq u(x_{1}), \dots, u_{n}:= u(x_{n}))\), a definição do
equivalente certo segue a forma
\[
u(c(L, u)) = \sum \limits_{k} p_{k} u_{k} = U(L)
\]

Pensando em termos de gráficos: \(c\) é o valor (não-aleatório,
``\textit{certo}'') com o qual o indivíduo alcança a utilidade média
(dadas as suas possibilidades \(u_{1},\dots, u_{n}\)). Intuitivamente,
pensaríamos que este ponto seria
\[
c = \sum \limits_{i=1}^{n} p_{i} x_{i}
\]
i.e.~o valor esperado dentre as possibilidades \(x_{i}\). Este caso,
porém, é o que se chama de ``indiferente ao risco''. Indivíduos que têm
\textit{aversão} ao risco, por outro lado, preferem um valor
\textit{menor, porém certo}, para evitar a aleatoriedade; eles se
dispõem a receber um valor menor, contanto que não se indisponham a
``recebíveis'' incertos. \textit{Mutatis mutandis} para indivíduos
\textit{amantes} ao risco: eles precisam ser pagos valores certos
\textit{maiores} para que se abstenham do risco.

!445 Fonte: \textcite{Cowell2004}, p.~191. Na figura, os possíveis
acontecimentos são \(x_{\color{red} Red}, x_{\color{blue} Blue}\)
(implicitamente com probabilidades respectivas
\(p_{\color{red} Red}, p_{\color{blue} Blue}\)), \(\mathcal{E}x\) é o
valor esperado dos acontecimentos \(x\) (i.e.~\(\sum p_{i}x_{i}\)), e
\(\xi\) é o equivalente certo para este indivíduo. Note-se que, como
\(\xi < \mathcal{E}x\), este indivíduo é \textit{avesso} ao risco; isso
também é notável por sua utilidade \textit{côncava}
(\(u^{\prime\prime}<0\)).
\end{definition}

\begin{definition}
Dado um indivíduo com um equivalente certo \(c(L, u)\) e cujos possíveis
acontecimentos \(x_{1},\dots,x_{n}\) (com resp. probabillidades
\(p_{1},\dots,p_{n}\)) tenham valor esperado
\[
\left<x\right> \coloneqq \sum \limits_{i=1}^{n} p_{i} x_{i}
\]
define-se seu \textit{prêmio de risco} como
\(\left< x \right> - c(L,u)\)
\parencite[, p. 191]{Cowell2004}.\footnote{Note-se que esta definição
  \textbf{não} é a mesma que \textcite{Mas-Colell1995} faz: MWG definem
  o que chamaram de ``\textit{probability premium}'' (definição
  \textbf{6.C.2 (ii)}, p.~186), mas não o \textit{risk premium}
  \textit{per se}.}

Para indivíduos indiferentes ao risco, ele é \(= 0\); para os avessos ao
risco, é \(>0\); para os amantes ao risco, é \(<0\).
\end{definition}

\hypertarget{exercuxedcios}{%
\section{Exercícios}\label{exercuxedcios}}

\begin{exercise}
Mostre que o equivalente certo de uma função utilidade exponencial
\(u(x) = - e^{-ax}\) assume a forma de uma \textit{média harmônica
ponderada}, ao se escrever \(x_i = \ln \frac{\tilde{x_{i}}}{a}\):
\[
C(L, a) = \frac{1}{a} \ln \left( \frac{1}{\sum \limits_{i} p_{i} \tilde{x}_{i}^{-1}} \right)
\]

\end{exercise}

\begin{exercise}
Mostre que o equivalente certo de uma função utilidade isoelástica
\(u(x) = \frac{x^{1-\rho}-1}{1-\rho}\) assume a forma de uma
\textit{\((1-\rho)\)-norma ponderada}, i.e.~
\[
   C(L, \rho) = \left( \sum \limits_{i} p_{i} x_{i}^{1-\rho} \right)^{1/(1-\rho)}
   \]

\end{exercise}

\begin{exercise}
Mostre que o equivalente certo de uma função utilidade logarítmica
\(u(x) = \ln x\) tem a forma de uma média geométrica ponderada
\[
C(L) = \prod \limits_{i=1}^{n} x_{i}^{p_{i}}
\]

\end{exercise}

\begin{exercise}
\parencite[, p. 135-6]{Mas-Colell1995} Seja \(F(x)\) uma função produção
com retornos de escala \textit{não-decrescentes}. Ou seja,
\[
\forall \alpha \in \mathbb{R}: F(\alpha x) \geq \alpha F(x)
\]
Prove que, dados preços \(p\), o seu lucro \(\pi(p)\) ou tende a
\(\infty\) ou \(\pi(p) \leq 0\).
\end{exercise}

\hypertarget{resoluuxe7uxf5es}{%
\section{Resoluções}\label{resoluuxe7uxf5es}}

\begin{exercise}
Dada uma loteria \(L = (p_{1},\dots,p_{n})\) de possíveis resultados
\(x=(x_{1},\dots,x_{n})\), o equivalente certo \(C(L, u)\) é a
quantidade monetária que garante a utilidade média (i.e.~utilidade de
von Neumann-Morgenstern). Como a utilidade \(u(x) = - e^{-ax}\) é
definida pelo coeficiente \(a\), podemos escrever \(C(L, u) = C(L, a)\).

Portanto,
\[
\begin{align*}
\cancel{ - }e^{-aC} &= \sum \limits_{i=1}^{n} p_{i} (\cancel{ - }e^{-ax_{i}}) \\
&= \sum \limits_{i=1}^{n} p_{i} e^{-ax_{i}}
\end{align*}
\]
Isolando \(C\), temos
\[
\begin{align*}
C(L, a) &= -\frac{1}{a} \ln\left( \sum \limits_{i=1}^{n}p_{i} e^{-ax_{i}} \right) \\
&= \frac{1}{a}\ln \left( \frac{1}{\left( \sum \limits_{i=1}^{n}p_{i} e^{-ax_{i}} \right)} \right)
\end{align*}
\]
Defina-se, a fins de simplificação, variáveis auxiliares
\(\tilde{x}_{i}\) tais que
\[
x_{i} = \frac{1}{a} \ln \tilde{x}_{i} \iff \tilde{x}_{i}^{-1} = e^{-ax_{i}}
\]
Dessa forma, temos que \$\$ \begin{align*}
C(L, a) &= \frac{1}{a}  \ln \left( \frac{1}{\left( \sum \limits_{i=1}^{n}p_{i} \frac{1}{\tilde{x}_{i}} \right)} \right) 

\end{align*} \$\$

Fazendo o mesmo para \(C\), i.e.~\(C = \frac{1}{a}\ln \tilde{C}\), temos
\$\$ \begin{align*}
\tilde{C} &= \frac{1}{\left( \sum \limits_{i=1}^{n}p_{i} \frac{1}{\tilde{x}_{i}} \right)}  \\
&= \frac{\sum \limits_{i=1}^{n} p_{i}}{\left( \sum \limits_{i=1}^{n}p_{i} \frac{1}{\tilde{x}_{i}} \right)} 

\end{align*} \$\$

Ou seja, é possível escrever o equivalente certo de uma utilidade
exponencial como uma média \textit{harmônica ponderada}.

A fins de \textit{sanity check}: suponhamos que algum dos \(x_{i}\)'s
tende a \(0\)\footnote{Em verdade: consideremos \(x_{i}\) se aproximando
  de \(0\). Não queremos tomar o limite formal
  \(\lim\limits_{x\to_{0}}\) neste caso.} --- \textit{mas não sua
probabilidade} \(p_i\) ---, ou seja, é um caso que tende à perda máxima
do indivíduo com uma probabilidade não-nula.

Teremos, portanto, que, neste caso extremo,
\[
C \to \frac{1}{a}  \ln \left( \frac{x_{i}}{p_{i}}\right) 
\]

Ou seja, o equivalente certo torna-se (muito mais) dependente deste caso
extremo: o quanto se desejará pagar para evitar a aleatoriedade desta
loteria será contingente
\end{exercise}

\begin{exercise}
Com a utilidade \(u(x) = \frac{x^{1-\rho}-1}{1-\rho}\), temos que
\(C(L, u) = C(L, \rho)\), obtido da seguinte forma:
\[
\begin{align*}
\frac{C^{1-\rho}-\cancel{ 1 }}{\cancel{ 1-\rho }} &= \sum \limits_{i=1}^{n} p_{i} \frac{x_{i}^{1-\rho}-\cancel{ 1 }}{\cancel{ 1-\rho }} \\
\end{align*}
\]

Isolando \(C\), temos
\[
C(L, \rho) = \left( \sum \limits_{i=1}^{n} p_{i} x_{i}^{1-\rho} \right)^{1/(1-\rho)}
\]

Ou seja, o equivalente certo de uma utilidade isoelástica é uma
\href{https://en.wikipedia.org/wiki/Norm_(mathematics)\#p-norm}{p-norma}\footnote{É
  uma \(p\)-norma \textit{de facto} assumindo-se que \(x_{i}>0\).} ---
no caso, uma \((1-\rho)\)-norma --- \textbf{ponderada} (!!!).
\end{exercise}

\begin{exercise}

\[
\begin{align*}
\ln C &= \sum \limits_{i=1}^{n} p_{i} \ln x_{i} \\
&= \sum \limits_{i=1}^{n} \ln x_{i}^{p_{i}} \\
&= \ln \left( \prod \limits_{i=1}^{n} x_{i}^{p_{i}} \right)
\end{align*}
\]

Portanto,
\[
\begin{align*}
C &= \prod \limits_{i=1}^{n} x_{i}^{p_{i}} \\
&= \left( \prod \limits_{i=1}^{n} x_{i}^{p_{i}} \right)^{1/\sum \limits_{i}p_{i}}
\end{align*}
\]

Ou seja, o equivalente certo de uma utilidade logarítmica é uma média
geométrica \textit{ponderada}. Note-se que, aqui, não é uma média
\textit{aritmética}, ou seja, estamos falando de variações
\textit{multiplicativas} de riqueza, muito propício p.~ex. em contextos
financeiros.
\end{exercise}

\begin{exercise}

\end{exercise}

\[
\phi\left( \frac{1}{t}z + \left( 1-\frac{1}{z} \right)z^{\prime} \right) 
\]

\printbibliography

\end{document}
